\section{Introduction}

Step-free rail access depends on both infrastructure provision and the continuity of the vertical links that make that infrastructure usable. A station can be classified as accessible while an elevator connecting the street, mezzanine, or platform is unavailable for part of the month. For travellers who cannot substitute stairs, elevator operation is therefore a basic service-side condition of step-free travel. This condition is especially visible in legacy rail systems, where accessible paths often depend on a limited number of vertical circulation assets.

Transport geography has long treated accessibility as a relationship between transport supply, spatial opportunities, and the constraints that shape travel \citep{geurs2004}. The updated accessibility framework also retains an explicit temporal component and recognises that different operationalisations can represent different aspects of accessibility \citep{geurs2026}. Recent work has shown how strongly measurement choices can shape spatial diagnosis. \citet{qu2025} found that the spatial resolution used to represent public transport access changes both the magnitude and geography of measurement error. \citet{schroth2026} similarly showed that an adjusted accessibility indicator can reveal over- and underserved areas that conventional measures do not identify. These studies make the construction of an indicator part of the analytical problem rather than a neutral technical step.

A related literature distinguishes scheduled transport supply from the service that is actually delivered. Real-time disruptions reduce realised accessibility relative to scheduled conditions \citep{liu2024}; recurring delays alter job-accessibility estimates \citep{nichols2024}; and travel-time variability can change measured accessibility inequalities \citep{braga2026}. For mobility-disabled travellers, operational reliability matters because the feasible network is already more constrained than for travellers able to use stairs or non-accessible alternatives \citep{liu2023}. Real-time elevator status has also been treated directly as an accessibility-information problem. \citet{sharif2022} collected elevator outage information across six North American rail systems and visualised whether nominally accessible stations were affected at a given time. That work demonstrates the practical importance of elevator operation. The question here is different: how does the definition of reliability change a long-run spatial diagnosis of elevator service?

New York City provides an unusually detailed case for examining that question. The Metropolitan Transportation Authority (MTA) publishes monthly equipment-level elevator availability over more than a decade. The published rate describes the share of reported service time during which an elevator is available. City accessibility reporting further notes that elevators undergoing planned replacement are excluded from the availability rate \citep{mopd2025}. This convention is well suited to evaluating equipment while it is counted as in service. It leaves a second monitoring question open: how much of the full calendar period is the asset recorded as available to provide elevator service?

That distinction matters geographically when excluded time is concentrated rather than evenly distributed. A systemwide availability rate can remain high while particular station complexes experience long periods outside the reported denominator. In that situation, changing the denominator can alter the relative ranking of places even when the published metric is internally consistent. Both indicators are internally valid; the question is whether they identify the same locations as persistently low-reliability.

This paper examines that problem using the New York City subway. Existing research has documented disability-related transit disparities \citep{liu2023}, realised accessibility losses under disruptions \citep{liu2024}, and additional travel burdens created by elevator placement within railway stations \citep{arai2026}. Open-data work has also represented New York subway vertical paths and elevator redundancy \citep{twosigma2020}. The present study stays at the asset and station-complex levels and isolates a prior measurement question: whether the temporal denominator used for elevator reliability changes the geography of long-run service-side monitoring.

Three research questions structure the analysis. First, how large and spatially variable is the gap between MTA-reported elevator availability and calendar service continuity? Second, does denominator choice change which station complexes are identified as persistently low-reliability locations, and are those relative rankings stable across years? Third, how closely does the asset-level geography correspond to the MTA's separate accessible-platform availability metric at borough scale?

These questions imply a simple spatial test. If denominator-excluded time is distributed proportionally across the subway system, calendar continuity should mainly shift the level of the indicator while preserving the ranking of places. If excluded time clusters in particular station complexes, the alternative denominator should change relative rankings and the set of locations repeatedly identified as low-reliability. The analysis therefore focuses on spatial reclassification and longitudinal rank persistence rather than the citywide mean alone.

The contribution is a measurement framework with a spatial consequence for transport geography. Calendar continuity complements MTA-reported availability by retaining calendar time that is absent from the operator's reliability denominator. Comparing the two reveals when an apparently technical denominator choice changes place-based monitoring. The resulting service geography is not a passenger-accessibility measure by itself. It is a service-side input that can be linked to path redundancy, passenger exposure, or population characteristics in subsequent accessibility and equity analyses.

\section{Operational reliability and step-free infrastructure}

\subsection{From nominal provision to service continuity}

Accessibility measures are selective representations of transport systems. Their usefulness depends on whether their assumptions match the decision problem being studied \citep{geurs2004}. In public transport, scheduled service provides a reproducible baseline, yet realised conditions can diverge substantially. \citet{liu2024} distinguished scheduled accessibility from accessibility under disruptions, while \citet{nichols2024} showed that recurring delays reduce realised access to employment in geographically uneven ways. \citet{braga2026} further demonstrated that incorporating travel-time variability can change the scale of measured inequality.

The same measurement logic applies to vertical access infrastructure. A binary indicator that a station contains an elevator describes provision, but it does not describe the time dimension of service. Transit performance guidance has long treated elevator reliability as an operational performance measure whose interpretation depends on service time and maintenance records \citep{tcrp2003}. An equipment availability percentage is therefore meaningful only in relation to the denominator used to define exposure. When long replacement or capital-work periods are excluded, the resulting percentage describes conditional in-service reliability. A calendar-time denominator answers a different question about continuity over the full month.

For step-free travel, the asset metric captures one service-side layer. Passenger accessibility also depends on path redundancy: one unavailable elevator may remove the only feasible vertical link at one station while remaining redundant at another. Research on mobility disabilities shows that formal transit provision and actual use can differ substantially \citep{liu2023}. Recent railway research likewise shows that elevator placement can impose additional travel-time burdens even at stations that are formally accessible \citep{arai2026}. Elevator service continuity therefore represents one layer of a broader accessibility system.

\subsection{Why the geography of reliability matters}

A systemwide average can conceal persistent local service loss. If particular station complexes repeatedly occupy the lower tail of the reliability distribution, the pattern is spatially structured rather than a sequence of isolated failures. That structure matters because a place that repeatedly experiences low continuity warrants different diagnosis from a systemwide random outage risk.

Recent JTG research increasingly treats transport reliability as a spatial problem. Recurring delays produce geographically differentiated accessibility losses \citep{nichols2024}; disruptions affect the reliability of delivered accessibility \citep{liu2024}; and service variability interacts with urban structure to change measured inequality \citep{braga2026}. Related work separates service-side signals from opportunity-side disadvantage because different indicators can identify different spatial regimes \citep{kang2026}. Elevator data allow that logic to be examined at the infrastructure layer: the question is whether the vertical assets that support step-free movement provide continuous service in the same places over time.

New York is particularly useful because the MTA provides a long monthly equipment panel together with a separate borough-level accessible-platform indicator. These sources allow the measurement gap to be examined at system, borough, and station-complex scales, while the platform indicator provides an independent higher-level comparison of spatial ordering.
